\begin{abstract}
Can explanations serve as a useful supervision signal for knowledge distillation in vision-language models?
We investigate this question in a plug-in replacement setting: a ViT student is trained to replace the visual module of a VLM, while the original language model remains fixed.
This setting is stricter than standalone embedding evaluation because the distilled representation must still support image-conditioned generation after it is inserted back into the multimodal pipeline.
We compare global feature reconstruction with explanation-guided token reconstruction, where Grad-CAM saliency from a frozen teacher probe weights the visual tokens that the student should match most strongly.
On Mini-ImageNet image-description prompts, three independent VLM judges compare anonymized teacher and student answers while observing the original image.
On the shared subset of 437 images for which every judge returned a valid preference for every variant, the teacher remains preferred overall; however, every explanation-guided objective improves over plain global reconstruction.
Per-judge trends and qualitative examples indicate that this gain is associated with preserving localized visual evidence and avoiding severe grounding failures, while errors requiring broader scene context remain.
These findings show that explanations provide a useful distillation signal, although the distilled visual module is not yet a faithful replacement for the teacher.
\end{abstract}
